\chapter{Технологическая часть}
\section{Средства разработки}
В качестве языка программирования был выбран Swift, так как он является языком со статической типизацией и является основным языком для разработки на операционную систему MacOS, на которой выполнялась лабораторная работа.
\section{Реализация алгоритмов} 


\lstinputlisting[caption={Поиск полным перебором}, label={lst:brute_force_search}]{brute_force_search.swift}
\lstinputlisting[caption={Нерекурсивный бинарный поиск}, label={lst:bin_search_1}, firstline=3, lastline=18]{bin_search_1.swift}
\lstinputlisting[caption={Нерекурсивный бинарный поиск с предварительным выходом}, label={lst:bin_search_2}, firstline=3, lastline=18]{bin_search_2.swift}
\lstinputlisting[caption={Рекурсивный бинарный поиск}, label={lst:bin_search_r}, firstline=5, lastline=20]{bin_search_recursion.swift}
\lstinputlisting[caption={Оберточная функция для рекурсивного бинарного поиска}, label={lst:wrapper}, lastline=3]{bin_search_recursion.swift}

Для единообразия интерфейсам функци при реализациия рекурсивного бинарно поиска \ref{lst:bin_search_r} используется оберточная функция \ref{lst:wrapper}

\section{Тестирование}
\begin{table}[H]
\caption{Функциональные тесты}
\label{tbl:tests}
\centering
\begin{tblr}{
  width = \linewidth,
  colspec = {Q[c,m] X[2,c,m] Q[c,m] Q[c,m] Q[c,m] X[1.5,c,m]},
  hlines, vlines,
  row{1,2} = {font=\bfseries},
}
\SetCell[r=2]{c} №
  & \SetCell[c=2]{c} Входные данные &
  & \SetCell[c=2]{c} Выходные данные &
  & \SetCell[r=2]{c} Комментарий \\
  & Массив & $x$ & Ожидаемые & Полученные & \\
1  & [1, 2, 3, 4, 5, 6, 7, 8]  & 1    & 0         & 0         & Первый элемент \\
2  & [1, 2, 3, 4, 5, 6, 7, 8]  & 8    & 7         & 7         & Последний элемент \\
3  & [1, 2, 3, 4, 5, 6, 7, 8]  & 5    & 4         & 4         & Середина, чётная длина \\
4  & [1, 2, 3, 4, 5, 6, 7]     & 4    & 3         & 3         & Середина, нечётная длина \\
5  & [1, 2, 3, 4, 5, 6, 7, 8]  & 10   & Не найден & Не найден & Больше максимума \\
6  & [3, 4, 5, 6, 7, 8, 9, 10] & 1    & Не найден & Не найден & Меньше минимума \\
7  & [1, 2, 3, 4, 7, 8, 9, 10] & 5    & Не найден & Не найден & Пропуск в середине \\
8  & [10, 20, 30, 40]          & 25   & Не найден & Не найден & Между соседними значениями \\
9  & [42]                      & 42   & 0         & 0         & Один элемент, найден \\
10 & [42]                      & 7    & Не найден & Не найден & Один элемент, не найден \\
11 & [1, 2]                    & 1    & 0         & 0         & Два элемента, первый \\
12 & [1, 2]                    & 2    & 1         & 1         & Два элемента, второй \\
13 & [1, 2]                    & 3    & Не найден & Не найден & Два элемента, не найден \\
14 & [$-10$, $-5$, $-1$, 0, 3] & $-5$ & 1         & 1         & Отрицательные значения \\
15 & [1, 100, 1000, 10\,000]   & 1000 & 2         & 2         & Разреженный массив \\
\end{tblr}
\end{table}
Все тесты пройдены успешно

\section{Вывод}
В результате технологической части были выбраны средства реализации,
реализованы алгоритм поиска полным перебором, два варианта нерекурсивного
бинарного поиска и рекурсивный бинарный поиск, а также проведено их
функциональное тестирование: все реализации прошли тесты, приведённые
в таблице~\ref{tbl:tests}.


